% Drawing 33/09-4 -- folder 33, batch 09, archive page 175 (PDF page 176).
% Transcription: un rectangle dont le coin supérieur droit est marqué $2d$ ; deux axes $(p, q)$, des points sur l'axe horizontal jusqu'à $2d$.
% Kind: rectangle; axes with points. Method 2 (schematic redraw from the facsimile). Status: ok.
% Notes: A rectangle whose upper right corner is a point marked 2d; to the right two axes (q upwards, p to the right), small points on the vertical near the origin and a point at 2d on the horizontal.
% Coordinates: 1 unit = 1 mm = 1/100 of the crop width (bbox in page fractions: [0.45, 0.6, 0.98, 0.75]).
% Redrawn by Claude (claude-opus-5-5), issue #73, 2026-11-07; not checked by a human; not compiled (no TeX in the session).
% Rights: reproduces part of the work -- behind RENDER_TRANSCRIPTS, never published by itself.
\documentclass[tikz,border=3pt]{standalone}
\usepackage{amsmath,amssymb}
\usetikzlibrary{patterns,arrows.meta}
\tikzset{pencil/.style={black!45}, struck/.style={black!35}, mark/.style={circle,draw,inner sep=1.3pt}, pt/.style={circle,fill,inner sep=0.9pt}, opt/.style={circle,draw,fill=white,inner sep=1.1pt}}
\begin{document}
\begin{tikzpicture}[x=1mm,y=1mm,line cap=round,line join=round,font=\scriptsize]
\draw (23,-15) -- (46,-15) -- (45,-33) -- (23,-33) -- cycle;
\node[pt] at (46,-15) {};
\node[inner sep=1pt] at (47,-12) {\tiny $2d$};
\draw (66,-13) -- (66,-33);
\draw (64,-32) -- (88,-31);
\node[pt] at (66,-32) {};
\node[pt] at (66,-30) {};
\node[pt] at (66,-28) {};
\node[pt] at (84,-31) {};
\node[inner sep=1pt] at (66,-11) {\tiny $q$};
\node[inner sep=1pt] at (90,-31) {\tiny $p$};
\node[inner sep=1pt] at (84,-34) {\tiny $2d$};
\end{tikzpicture}
\end{document}
